\section{Application to Gaussian Random  Design}\label{sec:example}

%In this section, we give two examples to illustrate that ReLU neural networks are adaptive to data distribution.
In literature, it is common to consider random Gaussian  design in nonparametric regression \citep{anderson1962introduction,muller2006linear,chatfield2018introduction}. In this section, we take anisotropic multivariate Gaussian design as example to justify Assumption \ref{assumption:pdata_cover} and demonstrate the effective Minkowski dimension. Here we only provide our main theorems and lemmas. The detailed proofs are given in Section \ref{sec:gaussian}. 

Consider a Gaussian distribution $P_{\rm data} \sim N(\bold{0},\Sigma)$ in $\RR^d$. The covariance matrix $\Sigma$ has the eigendecomposition form: $\Sigma =Q\Gamma Q^\top$, where $Q$ is an orthogonal matrix and $\Gamma = \diag(\gamma_1,\ldots,\gamma_d)$. For notational convenience in our analysis, we further denote eigenvalue $\gamma_i=\lambda_i^2$ for $i=1,\ldots,d$. Without loss of generality, assume that $\lambda_1^2\geq \lambda_2^2 \geq \ldots \geq \lambda_d^2$. Furthermore, we assume that $\{ \lambda_i^2 \}_{i=1}^d$ has an exponential or polynomial decay rate:
\begin{assumption}[Exponential decay rate]\label{assump:eigen-decay-exp}
    The eigenvalue series $\{\gamma_i \}_{i=1}^d=\{\lambda_i^2 \}_{i=1}^d$ satisfies $\lambda_i \leq \mu \exp\{-\theta i\}$ for some constants $\mu,\theta>0$.
\end{assumption}
\begin{assumption}[Polynomial decay rate]\label{assump:eigen-decay-poly}
    The eigenvalue series $\{\gamma_i \}_{i=1}^d=\{\lambda_i^2 \}_{i=1}^d$ satisfies $\lambda_i \leq \rho i^{-\omega}$ for some constants $\rho>0$ and $\omega > 1$.
\end{assumption}

When the eigenvalues decay fast, the support of the data distribution $P_{\rm data}$ has degeneracy in some directions. In this case, the majority of probability lies in some region $S \subset \RR^d$, which has an effective Minkowski dimension $p<d$. %it suffices to focus some $p$-dimensional subspace since the other $d-p$ dimensions degenerate in probability measure, where $p$ is the effective Minkowski dimension of data distribution. Therefore, we design $S$ in a way such that $P_{\rm data}(S)$ concentrates in  $p$ dimensions. %Intuitively, the number of hypercubes with side length $r$ used to cover $S$ is approximately $r^{-p}$.In literature, it is common to truncate the ambient space within a hyper-ellipsoid for Gaussian distribution. Similarly, we 
Specifically, consider a \lq\lq thick\rq\rq\ low-dimensional hyper-ellipsoid in $\RR^d$, 
\begin{align}
    S(R,r;p)\! :=\! \bigg\{ Qz \bigg| z\!=\!(z_1,\ldots,z_d)\in \RR^d, \sum_{i=1}^p \frac{z_i^2}{\lambda_i^2} \leq R^2, |z_j| \leq \frac{r}{2} \text{ for } j=p+1,\ldots,d  \bigg\}, \label{eq:set}
\end{align}
where $R,r>0$ and $p\in \NN_+$ are independent parameters. For the simplicity of notation, we first define a standard hyper-ellipsoid and then linearly transform it to align with the distribution $N(0,\Sigma)$.  The set $S(R,r;p)$ can be regarded as a hyper-ellipsoid scaled by $R>0$ in the first $p$ dimensions, and  with thickness $r>0$ in the rest $d-p$ dimensions. Then we construct a minimal cover as a union of nonoverlapping hypercubes with side length $r$ for $S(R,r;p)$.
The following lemma characterizes the relationship between the probability measure outside $S(R,r;p)$ and its covering number.

\begin{lemma}\label{lemma:gaussian}
    Given the eigenvalue series $\{\lambda_i^2\}_{i=1}^d$, for any $R,r>0$, choose $p>0$ such that $\lambda_p^{-1} = 2R/r$. If $p< R^2$, we will have
    \begin{equation*}
        \PP(X \notin S(R,r;p) ) =
        O \bigl(\exp( -R^2/3) \bigr),
    \end{equation*}
    \begin{equation*}
      N_r(S(R,r;p)) \leq \biggl(\frac{2R}{r} \biggr)^p \cdot\prod_{i=1}^p \lambda_i= \prod_{i=1}^p \bigg( \frac{\lambda_i}{\lambda_p}\bigg).
    \end{equation*}
    %where $N_r(S(R,r;p)$ denotes the covering number of $S(R,r;p)$ using hypercubes with side length $r$.
\end{lemma}
\begin{remark} 
Since data distribution $P_{\rm data}$ is supported on $\RR^d$, both the intrinsic dimension and the Minkowski dimension of $P_{\rm data}$ are $d$. However, Lemma \ref{lemma:gaussian} indicates that the effective Minkowski dimension of $P_{\rm data}$ is at most $p$. 
\end{remark}

According to Lemma \ref{lemma:gaussian}, if we choose scale $R>\sqrt{p}$ properly, the probability outside $S$ can be sufficiently small while the covering number of $S$ is dominated by $r^{-p}$, which gives that the effective Minkowski dimension of $P_{\rm data}$ is at most $p$. Moreover, under fast eigenvalue decays, the product of the first $p$ eigenvalues appearing in $N_r(S(R,r;p))$ is a small number dependent of $p$. In these cases, we specify the selection of $R$, $r$ and $p$ accordingly and show the effective Minkowski dimension is reduced to $p/2$ in Appendix \ref{sec:dim-gaussian}. 

Furthermore, we remark that the effective Minkowski dimension $p$  is not a fixed number given data distribution $P_{\rm data}$, but an increasing function of sample size $n$. As sample size $n$ increases, the estimation accuracy of $f^*$ is required to be higher, so that we are supposed to design more and smaller hypercubes to enable preciser estimation by neural networks. Besides, some of the $d-p$ dimensions are not negligible anymore and thereby become effective compared to the accuracy. Therefore, we need to incorporate more dimensions to be effective to achieve higher accuracy. 



With this observation, we construct $S(R,r;p)$ such that its effective Minkowski dimension $p(n)$ increases while thickness $r(n)$ decreases as sample size $n$ grows to enable preciser estimation. Then we develop the following sample complexity:

\begin{theorem}[Generalization error under fast eigenvalue decay]\label{thm:generalization-exp}
Under Assumption \ref{assump:function}, let $\hat{f}$ be the global minimizer of empirical loss given in \eqref{eq: obj} with function class 
$\mathcal{F}=\mathcal{F}(L,B,K)$. Suppose Assumption \ref{assump:eigen-decay-exp} hold. Set a tuple $(L,B,K)$ with the constants $C_1,C_2$ and $s$ appearing in Theorem \ref{lemma:approx} as 
\begin{equation*}
\begin{aligned}
    L= C_1,~~B = O\biggl(n^{\frac{\beta s}{2\beta+\sqrt{\log n/\theta}}} (\log n)^{\beta s} \biggr),~
    \text{and}~K = C_2 n^{\frac{\sqrt{\log n/\theta}}{2\beta + \sqrt{\log n/\theta} } }.
\end{aligned}
\end{equation*}
Then  we have 
\begin{equation*}
    \EE \norm{ \hat{f} - f^*}^2_{L^2(P_{\rm data})} = O\Bigl( \sigma^2 n^{- \frac{2\beta(1-\eta)}{2\beta+\sqrt{\log n/\theta}}}(\log n)^{3/2} \Bigr)
\end{equation*}
for sufficiently large $n$ satisfying $\log(\log n) / \sqrt{\theta \log n} \leq \eta$, where $\eta >0$ is an arbitrarily small constant.
Moreover, suppose Assumption \ref{assump:eigen-decay-poly} hold instead.  Set a tuple $(L,B,K)$ as \begin{equation*}
\begin{aligned}
        L= C_1,~~B&=O\Bigl(n^{\frac{(1+1/\omega)\beta s}{2\beta+ n^{\kappa}} } \Bigr),~\text{and}~K=C_2 \Bigl(n^{\frac{(1+1/\omega) n^{\kappa/(2\beta+n^\kappa)}}{4\beta + 2n^\kappa}} \Bigr),
\end{aligned}
\end{equation*}  
where $\kappa=(1+1/\omega)/\omega$. Then we have 
\begin{equation*}
     \EE \norm{ \hat{f} - f^*}^2_{L^2(P_{\rm data})} = O\Bigl( \sigma^2 n^{- \frac{2\beta}{2\beta+ n^{(1+1/\omega)/\omega}}} \log n\Bigr).
\end{equation*}
\end{theorem}

%\begin{theorem}[Generalization error under polynomial eigenvalue decay]\label{thm:generalization-poly}
%Fix any $f^*\in \mathcal{H}(\beta,[0,1]^d,M)$.  Set a tuple $(R,L,B,K)$ with the constants $C_1,C_2$ and $s$ appearing in Lemma \ref{lemma:approx} as $R=\|f^*\|_{\infty}$, $L=C_1$, $B=O(n^{2\beta s/(2\beta+p)})$ and $K=C_2n^{p/(2\beta+p)}$. Take the functional class 
%$\mathcal{F}=\mathcal{F}(R,L,B,K)$ in \eqref{eq: obj}. Then there exists a constant %$C=C(R,\beta,d,p,s,M,\sigma)$ such that 
%\begin{equation*}
%    \| \hat{f} - f^*\|^2_{L^2(P_{\rm data})} \leq C n^{-\frac{2\beta}{2\beta+n^{(1+\omega)/\omega^2}}} \log n.
%\end{equation*}
%holds with probability at least $1-2\exp(-n^{p/(2\beta+p)})$ for any $n \geq N$ with a sufficiently large $N$.
%\end{theorem}
Theorem \ref{thm:generalization-exp} suggests the effective Minkowski dimension of Gaussian distribution is $\sqrt{\log n/\theta}$ under exponential eigenvalue decay with speed $\theta$ and effective Minkowski dimension is $n^{(1+1/\omega)/\omega}$ under polynomial eigenvalue decay with speed $\omega$. For moderate sample size $n$, i.e. effective Minkowski dimension is less that ambient dimension $d$, Theorem \ref{thm:generalization-exp} achieves a faster convergence rate.  When we have a vast amount of data, the effective Minkowski dimension is the same as the ambient dimension $d$, and then we can apply standard analysis of deep neural networks for $d$-dimensional inputs to obtain the convergence rate  $\tilde{O}(n^{-2\beta/(2\beta+d)})$. To the best of our knowledge, Theorem \ref{thm:generalization-exp} appears to be the first result for nonparametric regression and deep learning theory, where the effective dimension varies with the sample size.

%\subsection{Application of Mixture Model}




